\section{Actual chart valuation witnesses in Shilov comparison}

\begin{definition}[Original chart test points]
\label{def:c0n-chart-test-points}
\uses{def:sg-div,lem:sg-div-points}
\lean{EnvelopeProperty.ProjSubvariety.IsChartTestPoint}
\leanok
A chart test point is either the original trivial point or a point
$v_{j,V,b}$ attached to an original chart index $j$ with $x_j\notin I_X$,
a genuine birational discrete valuation ring $V$ of
$B_j=A_j[t]$, and an integer $b>0$ for which
\[
 \operatorname{ord}_V\!\left(\sum_\ell f_\ell t^\ell\right)
 =\min_\ell\bigl(\operatorname{ord}_V(f_\ell)+\ell b\bigr).
\]
Here $V$ is essentially of finite type over the original field, has
fraction field $\operatorname{Frac}(B_j)$, and the analytic point is
exactly the original chart divisorial point of \cref{lem:sg-div-points}.
\end{definition}

\begin{lemma}[Actual chart orders evaluate whole polynomial generators]
\label{lem:c0n-chart-generator-orders}
\uses{def:c0n-chart-test-points,lem:sg-phi-basic,lem:sg-div-points}
\lean{EnvelopeProperty.ProjSubvariety.coneIdealPotential_chartTest_eq_finset_sup,EnvelopeProperty.ProjSubvariety.coneIdealPotential_chartTest_le_of_generator_orders,EnvelopeProperty.ProjSubvariety.chartTest_order_le_phiIdeal_obstacleIdeal}
\leanok
Fix actual chart data $(j,V,b)$ as in the definition.
Let $G$ be a finite set of whole ambient polynomials, each homogeneous of
one $x$-degree, and suppose that
$(\bar G)\subset B\subset(\bar G)^{\mathrm{sat}}$.
Then
\[
 \phi_B(v_{j,V,b})=
 \max_{F\in G}\left(-b^{-1}\operatorname{ord}_V(F^{(j)})\right).
\]
Consequently a common extended-real upper bound on these genuine generator
orders bounds the whole ideal potential.
For an original obstacle generator $F\in V_q(\phi)$, its chart order also
satisfies
$-b^{-1}\operatorname{ord}_V(F^{(j)})\le\phi_{J_q(\phi)}(v_{j,V,b})$.
\end{lemma}
\begin{proof}
The mapped-span and saturation bounds identify the potential of $B$ with
the potential of the ambient ideal generated by $G$ on $X^{\mathrm{an}}$.
Its potential is the finite maximum of the logarithmic values of the
whole polynomials in $G$.  The actual chart valuation identity identifies
each such logarithmic value with
$-b^{-1}\operatorname{ord}_V(F^{(j)})$, including the value $-\infty$
when the image is zero.  Taking the maximum proves the equality and its
upper-bound consequence.  Every original obstacle generator belongs to
the original obstacle ideal and has one homogeneous $x$-degree, so its
logarithmic value is at most that ideal potential.  Applying the same
chart valuation identity proves the last assertion.
\end{proof}

\begin{lemma}[An original generator attains the extended DVR ideal order]
\label{lem:c0n-chart-original-generator-order}
\uses{lem:c0n-original-member-attains-mapped-dvr-order}
\lean{EnvelopeProperty.Char0.exists_generator_eq_dvrIdealOrder_map_span}
\leanok
Let $\rho:R\to V$ be an original ring map into a DVR, and let $S\subset R$
have nonzero extended ideal $\rho_*((S))$.  Some original generator
$s\in S$ has order exactly $\operatorname{ord}_V(\rho_*((S)))$.
\end{lemma}
\begin{proof}
The order of the extended ideal is a finite natural number $r$ and is
attained by the image of an original element of $(S)$.
Every generator image has order at least $r$.  If no generator attained
$r$, all images of generators would have order at least $r+1$.
Their entire ideal would then lie in the $(r+1)$st maximal-ideal power,
including the image of the element of order $r$.  This would force
$r+1\le r$, a contradiction.
\end{proof}

\begin{lemma}[The original obstacle ideal controls its actual chart ideal order]
\label{lem:c0n-chart-obstacle-ideal-order}
\uses{lem:c0n-chart-original-generator-order,lem:c0n-chart-generator-orders,def:sg-Jm}
\lean{EnvelopeProperty.ProjSubvariety.chartTest_dvrIdealOrder_le_phiIdeal_obstacleIdeal}
\leanok
Let $I_q^{(j)}$ be the image of the entire original obstacle ideal $J_q(\phi)$
under the original chart dehomogenization map into $A_j[t]$.
If its extension to the actual chart DVR $V$ is nonzero, then
\[
 -b^{-1}\operatorname{ord}_V(I_q^{(j)}\mathcal O_V)
 \le \phi_{J_q(\phi)}(v_{j,V,b}).
\]
\end{lemma}
\begin{proof}
The obstacle ideal is generated by the original obstacle space
$V_q(\phi)$.  Apply the preceding generator-attainment lemma to the actual
composite from the original polynomial ring through chart
dehomogenization into $V$.  This gives an original obstacle generator
whose dehomogenized image attains the entire extended ideal order.
The last assertion of \cref{lem:c0n-chart-generator-orders} applied to
that generator proves the displayed inequality.
\end{proof}

\begin{lemma}[The Shilov boundary retains actual chart witnesses]
\label{lem:c0n-chart-shilov-core}
\uses{def:c0n-chart-test-points,lem:sg-finite-fs,lem:sg-rees-invariant,thm:bj-rees}
\lean{EnvelopeProperty.ProjSubvariety.exists_finset_chartTest_shilov_core,EnvelopeProperty.ProjSubvariety.exists_finset_chartTest_shilov_coneSup}
\leanok
Assume that $k$ is algebraically closed.  Let $s_i$ be finitely many forms
of one degree without a common zero on
$X^{\mathrm{an}}$, and let $e_i\in\mathbf N$.
There is a finite set $\Sigma$ of chart test points, containing the
original trivial point, such that for any polynomial $F$ whose
$t$-coefficients all have one homogeneous $x$-degree and any $r\in\mathbf N$,
\[
 |F|_{\sigma(w)}\le\max_i|s_it^{e_i}|_{\sigma(w)}^r
 \quad(w\in\Sigma)
\]
implies the same inequality at every $v\in X^{\mathrm{an}}$.
The equivalent logarithmic single-form comparison holds on the same set.
\end{lemma}
\begin{proof}
Take the union of the trivial point and all points attached to the actual
Rees valuation rings of the dehomogenized ideals
$(s_i^{(j)}t^{e_i})\subset B_j$, over the finitely many original charts.
This union is finite because each ideal has finitely many Rees valuations.
A power of $t$ belongs to each dehomogenized ideal.  Thus its Rees
valuations have positive $t$-order and satisfy the displayed torus law.
Their attached points are therefore chart test points with their actual
chart, valuation ring, and positive integer retained.
For an arbitrary $v$, choose a coordinate with $|x_j|_v=1$.
The Gauss seminorm on $B_j$ is at most one and agrees with
$|F|_{\sigma(v)}$ on dehomogenized homogeneous polynomials.
The assumed comparisons at all the Rees valuation points say that the
$r$th power of the dehomogenized ideal divides the dehomogenization of
$F$ in every Rees valuation ring.  The actual Rees valuation criterion
then gives the desired seminorm inequality at $v$.
For the logarithmic single-form version, the comparison at the trivial
point forces the integral weight $z$ to be nonpositive.  Write
$\log|s|+z=\log|st^{-z}|_{\sigma(\cdot)}$ and apply the polynomial version.
If $s$ vanishes on $X$, its logarithmic value is $-\infty$ and the
comparison holds directly.
\end{proof}

\begin{lemma}[Fubini--Study comparison on chart tests]
\label{lem:c0n-chart-fs-shilov}
\uses{lem:c0n-chart-shilov-core,lem:sg-finite-fs}
\lean{EnvelopeProperty.ProjSubvariety.IsFinFS.exists_finset_chartTest_le_of_le,EnvelopeProperty.ProjSubvariety.IsFinFS.exists_finset_chartTest_iSup_sub_eq}
\leanok
For every finite rational Fubini--Study function $\psi$ there is a finite
set $\Sigma$ of chart test points, containing the trivial point, such
that every rational Fubini--Study function bounded by $\psi$ on $\Sigma$
is bounded by $\psi$ everywhere.  On this same set the supremum of
$\varphi-\psi$ equals its finite maximum for every rational
Fubini--Study function $\varphi$ of any rational class.
\end{lemma}
\begin{proof}
Write $\psi=K^{-1}(\phi_{(s_it^{e_i})}+Z)$ using its finite normal form.
Clear the denominators of the weights of the competing Fubini--Study
function.  Each of its terms, after raising its section to the positive
integer power $KQ$, becomes one logarithmic single-form inequality
against an integral multiple of $\phi_{(s_it^{e_i})}$.  The preceding
lemma extends that inequality from $\Sigma$ to every point.
Taking the finite maximum of the terms gives the first assertion.
For the supremum assertion, choose a rational $q$ strictly between the
maximum on $\Sigma$ and any proposed larger value elsewhere.
The function $\varphi-q$ is bounded by $\psi$ on $\Sigma$, hence everywhere
by the first assertion.  Therefore $\varphi-\psi\le q$ everywhere,
contradicting that larger value.  The reverse inequality between the
finite maximum and the global supremum follows from inclusion.
\end{proof}

\begin{theorem}[Plurisubharmonic comparison needs only actual chart tests]
\label{thm:c0n-chart-test-comparison}
\source{bj22-trivval:thm:suppsh}
\uses{lem:c0n-chart-fs-shilov,thm:pg-psh2,lem:sg-pl-separate}
\lean{EnvelopeProperty.ProjSubvariety.IsPL.exists_finset_chartTest_iSup_sub_eq,EnvelopeProperty.ProjSubvariety.IsPL.exists_finset_chartTest_le_of_le,EnvelopeProperty.ProjSubvariety.IsPSH.le_of_forall_isChartTestPoint_le_isPL,EnvelopeProperty.ProjSubvariety.IsPSH.le_of_forall_isChartTestPoint_le}
\leanok
Let $k$ be algebraically closed.  For every piecewise linear $f$ there is
a finite set $\Sigma$ of chart test points containing the trivial point
such that
\[
 \sup_{X^{\mathrm{an}}}(u-f)=\max_\Sigma(u-f)
\]
for every plurisubharmonic $u$ of every real class.
Consequently $u\le h$ on all chart test points implies $u\le h$
everywhere whenever $h$ is upper semicontinuous and never $+\infty$.
\end{theorem}
\begin{proof}
Write $f=\phi_1-\psi_1$ with finite rational Fubini--Study functions and
choose the preceding finite set for $\phi_1$.
Approximate $u$ by a decreasing net $u_\alpha$ of finite rational
Fubini--Study functions.  Then
$u_\alpha-f=(u_\alpha+\psi_1)-\phi_1$, so its global supremum equals
its maximum over this same finite set.
If $u(v)-f(v)$ exceeded the finite maximum, choose a real number strictly
between them.  At every point of the finite set a sufficiently late
net member lies below that number; directedness supplies one index
working at all of them.  Its finite maximum then lies below the chosen
number, whereas its value at $v$ is at least $u(v)-f(v)$, a contradiction.
This proves the supremum formula and comparison with piecewise linear
functions.  Finally an upper semicontinuous function never equal to
$+\infty$ is the infimum of its strictly larger piecewise linear
functions.  Comparison with each of these majorants proves $u\le h$.
\end{proof}

\begin{theorem}[Polynomial envelope assembly using actual chart tests]
\label{thm:c0n-polynomial-chart-envelope-assembly}
\uses{thm:c0n-chart-test-comparison,lem:c0n-polynomial-cone-fs,lem:c0n-cone-ideal-bounds,thm:pg-psh2}
\lean{EnvelopeProperty.ProjSubvariety.isPSH_envelope_of_polynomialConeIdealChartEstimates}
\leanok
In \cref{thm:c0n-polynomial-cone-envelope-assembly}, it suffices to prove
the raw ideal discrepancy estimate at the original trivial point and at
all the actual chart points $v_{j,V,b}$ in
\cref{def:c0n-chart-test-points}.  The conclusion remains the
plurisubharmonicity of the original envelope on $X^{\mathrm{an}}$.
\end{theorem}
\begin{proof}
The continuous obstacle has a real lower bound $C$, and this constant is
a competitor for its envelope.  The ideal inclusions and bounded
polynomial generators give a decreasing factorial sequence of
Fubini--Study approximants, in classes tending to $d$, whose infimum $U$
dominates the envelope and hence $C$.  Thus $U$ is not identically
$-\infty$, and the decreasing-limit theorem makes $U$ plurisubharmonic.
The raw discrepancy estimates give
$U(v)\le P_d(\phi)(v)+A_v/j!$ at every chart test point.
Letting $j$ tend to infinity yields $U\le\phi$ there.
The preceding comparison theorem with the piecewise linear obstacle
therefore gives $U\le\phi$ everywhere.  Thus $U$ is a competitor for the
envelope and equals it.
\end{proof}
